\section{Introduction}
Il existe plusieurs méthodes pour caractériser les propriétés thermiques et mécaniques d’un solide. Dans le cadre de cette étude, nous analysons les propriétés mécaniques d’un échantillon de PMMA, et plus précisément son comportement en flexion (voir Fig. \ref{fig:meca}), lors d’une première expérience.
\begin{figure}[ht]
\centering
\includegraphics[scale=0.2]{meca.png}
\caption{Techniques d’analyse des propriétés mécaniques d’un solide, réf. : \cite{meca}}
\label{fig:meca}
\end{figure}
\\
En effet, nous nous concentrons, dans ce rapport, à l’étude du comportement des propriétés mécaniques des polymères en fonction de différentes rampes thermiques. Pour ce faire, nous utilisons des outils d’analyse thermique pour étudier l’évolution dynamique des propriétés des matériaux. En effet, pour les polymères, ces méthodes constituent un outil essentiel pour comprendre leur structure, leurs transitions physiques, leurs historiques thermiques et prédire leur comportement en service (Ref. \cite{protocoleDSC}, p. 2 ; Ref. \cite{protocoleDMA}, p. 2).
\\\\
Nous étudions, dans ce rapport, principalement deux techniques. La DMA (Dynamic Mechanical Analysis) mesure la réponse mécanique viscoélastique d’un polymère soumis à une contrainte sinusoïdale (Ref. \cite{protocoleDMA}, p. 2). La DSC (Differential Scanning Calorimetry), en revanche, mesure les flux de chaleur nécessaires pour maintenir l’échantillon et une référence à la même température lors d’une rampe thermique (Ref. \cite{protocoleDSC}, pages 2–4).
\\\\
Ainsi, le but du laboratoire est double. Premièrement, nous utilisons la DMA en configuration trois points pour mesurer les modules viscoélastiques $E'$, $E''$ et $\tan\delta$ (décrits ultérieurement, Sec. \ref{sec:theorie}) sur plusieurs fréquences et extraire l’énergie d’activation de la transition vitreuse (transition alpha, Ref. \cite{DMA}) de notre échantillon de PMMA. En second lieu, nous utilisons la DSC pour identifier les transitions thermiques (cristallisation froide, fusion, transition vitreuse, cristallisation chaude), les chaleurs latentes et le taux de cristallinité pour un échantillon de PET (Ref. \cite{DSC}).

%Il existe plusieurs manieres de verifier les proprietes mecanique d'un solide (voir Fig. \ref{fig:meca}). Dans notre cas, c'est par la flexion que nous analysons un echantillon de PMMA pour une premiere exeperience. 

%\begin{figure}[ht]
%\centering
%	\centering
 %       \includegraphics[scale=0.28]{meca.png}
  %      \captionof{figure}{ Techniques d'analyses des proporietes mecaniques d'un solide, ref : \cite{meca} }
%\label{fig:meca}
%\end{figure}
\\
%En effet, nous nous concentrons dans ce rapport a l'etude du comportement des proprietes mecaniques des polymeres en fonction de differentes rampes thermiques. Pour ce faire, nous utilisons des outils d'analyse thermique pour etudier l'evolution dynamique des proprietes des materiaux. Car pour les polymères, ces méthodes constituent un outil essentiel pour comprendre leur structure, leurs transitions physiques, leurs historiques thermiques et prédire leur comportement en service (Ref. \cite{protocoleDSC}, p.2 ; Ref. \cite{protocoleDMA}, p.2).
\\\\
%Nous etudions dans ce rapport principalement deux techniques. La DMA (Dynamic Mechanical Analysis)  mesure la réponse mécanique viscoélastique d’un polymère soumis à une contrainte sinusoïdale (Ref. \cite{protocoleDMA} p.2). La DSC qui en revanche (Differential Scanning Calorimetry) mesure les flux de chaleur nécessaires pour maintenir l’échantillon et une référence à la même température lors d’une rampe thermique (Ref. \cite{protocoleDSC} pages 2-4).
\\\\
%Ainsi le but du laboratoire est double. Premierement, nous utilisons la DMA en une configuration en trois points pour mesurer les modules viscoélastiques E’, E’’ et tan δ (decrits posterieurement Sec. \ref{sec:theorie}) sur plusieurs fréquences et extraire l’énergie d’activation de la transition vitreuse (transition alpha) Ref. \cite{DMA} de notre echantillon de PMMA. En un second temps, nous utilisons la DSC pour identifier les transitions thermiques (cristallisation froide, fusion, transition vitreuse, cristallisation chaude), les chaleurs latentes et le taux de cristallinité pour un echantillon de PET Ref.\cite{DSC}.
